\subsection{Beispiele}
\coloredbox{red}{REWRITE}{
    This section is undergoing a rewrite, to be more on-point. The previous section may also see some changes as a result.
}
\howto{Tests}{
    Was zu tun / erwähnen ist jeweils:
    \begin{enumerate}
        \item Modell: Verteilung der Stichproben $\cX_k$ unter $\P_\vartheta$, mit $\vartheta =$ gesuchte Werte.
              Für $\cX_k \sim \cN(\mu, \sigma^2)$, ist $\vartheta = (\mu, \sigma^2)$, wenn beide unbekannt.
        \item Hypothesen: $H_0: \vartheta = \vartheta_0 =$ Vergleichswert, $H_A: \vartheta < \vartheta_0$ (oder grösser, $\neq$, etc). \hl{Umgekehrt, s. links}
        \item Teststatistik: Basierend auf Verteilung von $\cX_k$ und Bekanntheit von Parametern
        \item Verteilung von $T$: Oft $t_{n - 1}$, siehe \ref{sec:estimators-properties} unten
        \item Verwerfungsbereich: Nach Anleitung unten
        \item Wert der $T$: $T(\omega)$ berechnen, $\omega$ ist beobachteter Wert (bspw. Durchschnitt $\overline{\cX}_n$)
        \item Testentscheid: Basierend auf Resultat von vorherigem und $K$ entscheiden.
    \end{enumerate}
}
\bi{Verwerfungsbereich} $K_l = (c_>, \8)$ (links), $K_r = (-\8, c_<)$ (rechts), $K_b = (-\8, c_{\neq}) \cup (c_{\neq}, \8)$ (beidseitig)

\shortexample[Gauss-Test] Voraussetz.:
\bi{(1)} $\cX_i$ i.i.d. $\sim \cN(\vartheta, \sigma^2)$ unter $\P_\vartheta$
\bi{(2)} $\sigma^2$ bekannt.
\textbf{Teststatistik}:
\[
    T = \frac{\overline{\cX}_n - \vartheta_0}{\frac{\sigma}{\sqrt{n}}} \sim \cN(0, 1) \qquad \text{unter } \P_{\vartheta_0}
\]
\bi{Bestimmung $c$s}: $z_{1 - \alpha} = \Phi^{-1}(1 - \alpha)$.\\
$c_> = z_{1 - \alpha}; c_< = -z_{1 - \alpha}; c_{\neq} = z_{1 - 0.5 \alpha}$


\shortexample[t-Test] $\cX_i$ i.i.d. $\sim \cN(\mu, \sigma^2)$ unter $\P_{\overrightarrow{\vartheta}}$
mit $\overrightarrow{\vartheta} := (\mu, \sigma^2)$ (und insbesondere $\sigma^2)$ unbekannt.
Testen wieder $H_0 : \mu = \mu_0$. Teststatistik ist hier (see \ref{sec:estimators-properties}, $t_m$-V.):
\[
    T = \frac{\overline{X}_n - \mu_0}{\frac{S}{\sqrt{n}}} \qquad S^2 = \frac{1}{n - 1} \sum_{k = 1}^{n} (\cX_k - \overline{\cX}_n)^2
\]
$K$: $c_> = t_{n - 1, 1 - \alpha}$; $c_< = t_{n - 1, \alpha} = -c_>$; $c_{\neq} = t_{n - 1, 1 - 0.5 \alpha}$,\\
mit $t_{m, \gamma}$ das $\gamma$-Quantil von $t_m$-Verteilung, welches Wert ist, so dass $\P[\cX \leq t_{m, \gamma}] = \gamma$ für $\cX \sim t_m$.
Wert ist tabelliert!


\shortremark[Zweistichproben-Tests] $\cX_k$ (i.i.d.) und jeweils unabhängig von allen $\cY_l$


\shortexample[Gepaarter Zweist. m. N.V.] $\cX_i$ i.i.d. $\sim \cN(\mu_\cX, \sigma^2)$ und $\cY_j$ i.i.d. $\sim \cN(\mu_\cY, \sigma^2)$.

Differenzen $\cZ_k = \cX_k - \cY_k$ unter $\P_\vartheta$ i.i.d. $\sim \cN(\mu_\cX, \mu_\cY, 2\sigma^2)$.
Wenn nicht unabhängig: $\sim \cN(\mu_\cX - \mu_\cY, 2(1 - \rho) \sigma^2)$ mit $\rho \in (-1, 1)$ Korrelation und $\rho = 0$ Unabhängigkeit.


\shortexample[Ungepaart] $\cX_i$, $\cY_j$ wie zuvor ($n$ und $m$ je).

\bi{$z$-Test} (falls $\sigma^2$ bekannt):
\[
    T = \frac{(\overline{\cX}_n - \overline{\cY}_m) - (\mu_\cX - \mu_\cY)}{\sigma \sqrt{\frac{1}{n} + \frac{1}{m}}} \sim \cN(0, 1)
\]
$\mu_\cX - \mu_\cY$ muss aus $H_0$ bekannt sein. Der kritische Bereich ist wieder Quantile der $\cN(0, 1)$-Verteilung.

\bi{$t$-Test} (falls $\sigma^2$ unbekannt):
\[
    T = \frac{(\overline{\cX}_n - \overline{\cY}_m) - (\mu_\cX - \mu_\cY)}{S \sqrt{\frac{1}{n} + \frac{1}{m}}} \sim t_{n + m - 2}
\]
mit $S^2_\cY$ analog zu $S^2_\cX$ ($n$ durch $m$ ersetzt)
\[
    S^2_\cX = \frac{1}{n - 1} \sum_{k = 1}^{n} (\cX_k - \overline{\cX}_n)^2
\]
Dann:
\[
    S^2 = \frac{1}{m + n - 2}((n - 1)S_\cX^2 + (m - 1)S_\cY^2)
\]
